\documentclass[a4paper,12pt]{article}
\usepackage[utf8]{inputenc}
\usepackage[T1]{fontenc}
\usepackage[french]{babel}
\usepackage{graphicx}
\usepackage{hyperref}
\usepackage{geometry}
\geometry{margin=2.5cm}
\usepackage{enumitem}
\usepackage{xcolor}
\usepackage{titlesec}
\usepackage{listings}
\lstset{basicstyle=\ttfamily\small, breaklines=true}

\title{\textbf{Cahier des charges technique} \\
       \Large Plateforme LocalGuide – Annuaire de commerces de proximité}
\author{Équipe projet LocalGuide}
\date{Version 1.0 – Mars 2026}

\begin{document}

\maketitle
\tableofcontents
\newpage

\section{Introduction}
LocalGuide est une plateforme web destinée à référencer, rechercher et localiser les commerces de proximité (cafés, restaurants, pharmacies, boutiques, services, etc.) sur le territoire de Monastir (Tunisie).  
L’objectif est double :
\begin{itemize}
    \item \textbf{Pour les citoyens / visiteurs} : découvrir facilement les commerces locaux, obtenir des informations pratiques (horaires, avis, photos) et planifier des parcours thématiques.
    \item \textbf{Pour les commerçants} : bénéficier d’une visibilité accrue et communiquer sur leurs actualités (promotions, événements).
\end{itemize}
Le site se veut moderne, intuitif et adapté à tous les supports (ordinateur, tablette, mobile).

\section{Public cible}
\begin{itemize}
    \item \textbf{Visiteurs} : habitants, touristes, personnes de passage souhaitant explorer les commerces de Monastir.
    \item \textbf{Commerçants} : propriétaires ou gérants de commerces locaux désireux d’ajouter et de gérer leur fiche.
    \item \textbf{Administrateurs} : modérateurs en charge de valider les contenus, de gérer les utilisateurs et de garantir la qualité des informations.
\end{itemize}

\section{Fonctionnalités principales}
D’après le document initial et les pages développées, les fonctionnalités suivantes sont implémentées :

\subsection{Annuaire des commerces}
Chaque commerce dispose d’une fiche détaillée contenant :
\begin{itemize}
    \item Nom, description, catégorie (restaurant, santé, service…).
    \item Adresse, horaires d’ouverture, téléphone, email.
    \item Photos, galerie d’images.
    \item Avis et notes des utilisateurs.
\end{itemize}

\subsection{Carte interactive (page d’accueil)}
La page principale affiche une carte de Monastir avec des marqueurs pour chaque commerce.  
Fonctionnalités :
\begin{itemize}
    \item Géolocalisation de l’utilisateur.
    \item Filtres par catégorie, distance, note.
    \item Clic sur un marqueur → ouverture de la fiche du commerce.
\end{itemize}

\subsection{Fiche commerce (page détail)}
Page dédiée à un commerce, avec :
\begin{itemize}
    \item Toutes les informations générales.
    \item Carte intégrée avec position exacte.
    \item Avis des visiteurs.
    \item Galerie photos.
    \item Boutons d’action (appeler, itinéraire, ajouter un avis).
    \item Commerces similaires suggérés.
\end{itemize}

\subsection{Ajout d’un commerce (formulaire)}
Formulaire permettant aux commerçants ou utilisateurs de proposer un nouveau commerce.  
Champs obligatoires : nom, catégorie, adresse, ville, code postal, téléphone, email, horaires.  
Optionnels : description, photo.  
Le commerce n’est pas publié immédiatement : il passe en \textbf{attente de validation} par un administrateur.

\subsection{Circuits / itinéraires (page guides)}
Création de circuits thématiques (gastronomique, shopping, touristique) pour faciliter la découverte de la ville.  
Chaque guide (personne référente) propose un parcours avec plusieurs commerces, affiché sur la carte dans un ordre logique.

\subsection{Actualités locales (page News)}
Espace dédié aux actualités des commerces et de la ville :
\begin{itemize}
    \item Promotions, événements, nouvelles ouvertures.
    \item Articles rédigés par l’équipe ou les commerçants.
    \item Possibilité de filtrer par type (actualité, événement, nouveau commerce, promo).
    \item Bandeau à la une pour l’événement principal.
    \item Inscription à une newsletter (simulation).
\end{itemize}

\subsection{Modération et administration}
Tableau de bord administrateur avec :
\begin{itemize}
    \item Liste des commerces en attente, approuvés, refusés.
    \item Filtres par statut.
    \item Actions : valider, refuser, voir le détail.
    \item Alertes (signalements, attente prolongée, etc.).
    \item Gestion des avis (suppression des abus).
    \item Gestion des utilisateurs.
\end{itemize}

\section{Description détaillée des pages réalisées}
\subsection{Accueil (accueil.html)}
\begin{itemize}
    \item Barre de navigation commune à toutes les pages.
    \item Diaporama d’images en arrière‑plan.
    \item Formulaire de recherche textuelle.
    \item Mots‑clés (tags) cliquables.
    \item Bandeau défilant (ticker) avec des informations dynamiques.
    \item Sélection de commerces en vedette (cartes avec note, catégorie, tags).
    \item Pied de page avec liens utiles.
\end{itemize}

\subsection{Ajout d’un commerce (add.html)}
\begin{itemize}
    \item Formulaire structuré en sections (informations générales, coordonnées, horaires et médias).
    \item Barre de progression visuelle.
    \item Validation des champs obligatoires avec messages d’erreur.
    \item Compteur de caractères pour la description.
    \item Zone de dépôt de photo (drag \& drop, sélection fichier).
    \item Boutons Annuler (confirmation) et Envoyer.
    \item Simulation d’envoi avec spinner et message de succès.
\end{itemize}

\subsection{Fiche commerce (fichecomm.html)}
\begin{itemize}
    \item Fil d’Ariane.
    \item En‑tête avec nom, catégorie, note moyenne, statut d’ouverture, adresse, horaires condensés.
    \item Boutons d’action (appeler, itinéraire, ajouter avis).
    \item Image de couverture.
    \item Grille à deux colonnes :
          \begin{itemize}
              \item Gauche : description, galerie photos, avis récents.
              \item Droite : informations de contact (téléphone, adresse, site, WhatsApp), carte statique (placeholder), horaires détaillés.
          \end{itemize}
    \item Deux exemples de commerces (Café Marina, Restaurant El Medina) pour illustrer la variété.
\end{itemize}

\subsection{Page Guides (guide.html)}
\begin{itemize}
    \item Présentation des guides locaux (personnes).
    \item Filtres par spécialité (Gastronomie, Histoire, Mer \& Port, Artisanat, Nocturne).
    \item Barre de recherche.
    \item Tri par note, nom, prix, nombre d’avis.
    \item Cartes guide avec photo de couverture, avatar, spécialité, description courte, tags, note, prix, disponibilité.
    \item Boutons de réservation (simulation) et de consultation de profil.
\end{itemize}

\subsection{Page Nouveautés (news.html)}
\begin{itemize}
    \item Onglets pour filtrer : Tout, Actualités, Événements, Nouveaux Commerces, Promos.
    \item Bannière « À la une » pour un événement majeur.
    \item Grille d’articles (actualités) avec image, catégorie, date, titre, extrait, auteur.
    \item Liste d’événements sous forme de carte (date, titre, lieu, heure, tarif, tag).
    \item Section d’inscription à la newsletter.
\end{itemize}

\subsection{Page Admin (admin.html)}
\begin{itemize}
    \item En‑tête avec badge Super Admin.
    \item Alertes de modération (signalements, attente prolongée).
    \item Filtres par statut (Tous, En attente, Approuvés, Refusés) avec compteurs dynamiques.
    \item Liste des commerces avec nom, adresse, horaires, statut et actions (Voir, Valider, Refuser).
    \item Pagination.
    \item Simulation des actions de validation / refus.
\end{itemize}

\section{Types d’utilisateurs et droits}
\begin{itemize}
    \item \textbf{Visiteur (non connecté)} : consultation de toutes les pages, recherche, lecture des avis, utilisation de la carte.
    \item \textbf{Commerçant} : peut soumettre un formulaire d’ajout ; après validation, pourra (dans une version future) gérer sa fiche et répondre aux avis.
    \item \textbf{Administrateur} : accès au tableau de bord de modération, validation des commerces, gestion des avis et utilisateurs.
\end{itemize}

\section{Contraintes techniques et recommandations}
\begin{itemize}
    \item \textbf{Développement} : HTML5, CSS3, JavaScript (vanilla). Les fichiers CSS sont séparés (style1.css, style\_nav\_barre.css, etc.).
    \item \textbf{Responsive design} : les pages doivent s’adapter aux mobiles et tablettes (utilisation de flexbox, grid, media queries).
    \item \textbf{Carte interactive} : à implémenter avec une bibliothèque comme Leaflet ou Mapbox (actuellement placeholders).
    \item \textbf{Accessibilité} : contrastes, labels de formulaire, alternative textuelle pour les images.
    \item \textbf{Sécurité} : validation des données côté client (JavaScript) et à prévoir côté serveur ; modération obligatoire avant publication.
    \item \textbf{Performances} : optimiser les images, minifier CSS/JS.
\end{itemize}

\section{Évolutions possibles}
\begin{itemize}
    \item Compte utilisateur avec authentification.
    \item Gestion complète des avis (ajout, modification, signalement).
    \item Système de réservation en ligne pour les circuits.
    \item Backend avec base de données (MySQL/PostgreSQL) et API REST.
    \item Internationalisation (arabe, anglais).
\end{itemize}

\section{Annexes – Structure des fichiers}
\begin{lstlisting}
Projet_web_LocalGuide/
│
├── accueil.html
├── add.html
├── admin.html
├── fichecomm.html
├── guide.html
├── news.html
├── style1.css
├── style_add.css
├── style_admin.css
├── style_fiche_comm.css
├── style_guide.css
├── style_news.css
├── style_nav_barre.css
└── images/ (logo, photos de demonstration)
\end{lstlisting}

\vfill
\hrule
\centering
\textit{Document rédigé à partir du projet initial et des pages web développées.}
\end{document}